\section{Auxiliary Lemmas}
\label{sec:basic-descent}
This section establishes the basic descent properties of the DDLALM algorithm.

For convenience, we introduce the successive differences
\[
    \Delta_{\bx}^{k+1}:=\bx^{k+1}-\bx^k,
    \qquad
    \Delta_{\by}^{k+1}:=\by^{k+1}-\by^k,
    \qquad
    \Delta_{\bmu}^{k+1}:=\bmu^{k+1}-\bmu^k.
\]
and set $\bar\rho_k:=\rho+s_k$.

Define
\[
    M_\mu:=\max\{\|\bmu^0\|,L_h\},
    \qquad
    C_r:=L_h+M_\mu,
\]
and introduce the splitting residual and the KKT-aligned multiplier
\[
\br^{k+1}:=\calA(\bx^{k+1})-\by^{k+1},
    \qquad
    \bz^{k+1}:=\bmu^k+\bar\rho_k\br^{k+1}.
\]

Subtracting $\bmu^k$ from \eqref{eq:mu-update} and using
\eqref{eq:lambda-update} gives the identity
\begin{equation}
    \Delta_{\bmu}^{k+1}
    =
    \tau_k\br^{k+1},
    \qquad k\ge0.
    \label{eq:mu-drift-tau}
\end{equation}

\begin{lemma}[Multiplier and splitting-residual bounds]
\label{lem:multiplier-boundedness}
\blue{Suppose that Assumption~\ref{ass:h} holds.}
Let $\{(\bx^k,\by^k,\bmu^k)\}_{k\ge0}$ be generated by
Algorithm~\ref{alg:ddlalm}.  Then, for
every $k\ge0$,
\begin{equation}
    \|\bmu^k\|\le M_\mu,
    \qquad
    \|\br^{k+1}\|
    \le
    \frac{C_r}{\bar\rho_k},
    \qquad
    \|\Delta_{\bmu}^{k+1}\|
    \le
    \frac{\tau_k C_r}{\bar\rho_k}.
    \label{eq:basic-bounds}
\end{equation}
\end{lemma}

\begin{proof}
The first-order optimality condition of the exact $\by$-update is
\[
    \zero
    \in
    \partial h(\by^{k+1})-\bmu^k-\bar\rho_k\br^{k+1},
\]
which proves $\bz^{k+1}\in\partial h(\by^{k+1})$.  \blue{By
\eqref{eq:bounded-h-subgradients}, $\|\bz^{k+1}\|\le L_h$.}

The identity
$\br^{k+1}=(\bz^{k+1}-\bmu^k)/\bar\rho_k$ and
\eqref{eq:mu-drift-tau} yield
\begin{equation}
    \bmu^{k+1}
    =
    (1-q_k)\bmu^k+q_k\bz^{k+1},
    \qquad
    q_k:=\frac{\tau_k}{\bar\rho_k}\in(0,1),
    \label{eq:mu-convex-combination}
\end{equation}
which is a convex combination.  Induction gives
$\|\bmu^k\|\le M_\mu$ for every $k\ge0$.

The definition of $\br^{k+1}$ now gives
\[
    \|\br^{k+1}\|
    \le
    \frac{\|\bz^{k+1}\|+\|\bmu^k\|}{\bar\rho_k}
    \le
    \frac{C_r}{\bar\rho_k},
\]
which proves the residual estimate in \eqref{eq:basic-bounds}.  Combining this
estimate with \eqref{eq:mu-drift-tau} gives the increment estimate, completing the proof.
\end{proof}

The lemma below characterizes the smoothness of
$\Phi_s(\bx,\by;\bmu)$ with respect to $\bx$ for fixed $\bmu$ and $\by$.

\begin{lemma}
\label{lem:x-block-smoothness}
\color{blue}
Suppose that Assumption~\ref{ass:smoothness} and the conditions of
Lemma~\ref{lem:multiplier-boundedness} hold.  Define
\[
    B_y
    :=\max\left\{
        \|\by^0\|,
        B_A+\frac{M_\mu+L_h}{\bar\rho_0}
    \right\}.
\]
Then $\|\by^k\|\le B_y$ for every $k\ge0$.  Moreover, for every $k\ge0$
and every $\bx,\bx'\in\calS$,
\begin{equation}
\begin{aligned}
&\left\|
\nabla_{\bx}\Phi_{s_k}(\bx,\by^k;\bmu^k)
-\nabla_{\bx}\Phi_{s_k}(\bx',\by^k;\bmu^k)
\right\|
\le L_k\|\bx-\bx'\|,\\
&\qquad
L_k
:=L_f+L_AM_\mu
+\bar\rho_k\left[L_A(B_A+B_y)+G_A^2\right].
\end{aligned}
\label{eq:x-block-smoothness-constant}
\end{equation}
In particular,
\begin{equation}
    L_k\le\widehat L\bar\rho_k,
    \qquad
    \widehat L
    :=\frac{L_f+L_AM_\mu}{\bar\rho_0}
      +L_A(B_A+B_y)+G_A^2.
    \label{eq:uniform-x-block-smoothness}
\end{equation}
\end{lemma}

\begin{proof}
\color{blue}
Set
\[
    \bu^{k+1}
    :=\calA(\bx^{k+1})+\frac{\bmu^k}{\bar\rho_k}.
\]
The optimality condition of the proximal update \eqref{eq:y-update} gives
\[
    \bar\rho_k(\bu^{k+1}-\by^{k+1})
    \in\partial h(\by^{k+1}).
\]
Hence \eqref{eq:bounded-h-subgradients}, the multiplier bound, and
$\bx^{k+1}\in\calS$ imply
\[
\begin{aligned}
    \|\by^{k+1}\|
    &\le \|\calA(\bx^{k+1})\|
      +\frac{\|\bmu^k\|+L_h}{\bar\rho_k}\\
    &\le B_A+\frac{M_\mu+L_h}{\bar\rho_0}
    \le B_y,
\end{aligned}
\]
where the second inequality uses the monotonicity of $\{\bar\rho_k\}$.
The definition of $B_y$ also covers $\by^0$.

Fix $k\ge0$ and $\bx,\bx'\in\calS$.  From \eqref{eq:reduced-al},
\[
    \nabla_{\bx}\Phi_{s_k}(\bx,\by^k;\bmu^k)
    =\nabla f(\bx)
     +J_{\calA}(\bx)^\top(\bmu^k-\bar\rho_k\by^k)
     +\bar\rho_kJ_{\calA}(\bx)^\top\calA(\bx).
\]
Because $\calS$ is convex, the Jacobian bound integrated along the segment
joining $\bx'$ and $\bx$ gives
\[
    \|\calA(\bx)-\calA(\bx')\|
    \le G_A\|\bx-\bx'\|.
\]
Moreover,
\[
\begin{aligned}
&J_{\calA}(\bx)^\top\calA(\bx)
-J_{\calA}(\bx')^\top\calA(\bx')\\
={}&\bigl(J_{\calA}(\bx)-J_{\calA}(\bx')\bigr)^\top\calA(\bx)
+J_{\calA}(\bx')^\top\bigl(\calA(\bx)-\calA(\bx')\bigr),
\end{aligned}
\]
and therefore
\[
\left\|
J_{\calA}(\bx)^\top\calA(\bx)
-J_{\calA}(\bx')^\top\calA(\bx')
\right\|
\le (L_AB_A+G_A^2)\|\bx-\bx'\|.
\]
Combining these estimates with
$\|\bmu^k-\bar\rho_k\by^k\|
\le M_\mu+\bar\rho_kB_y$
proves \eqref{eq:x-block-smoothness-constant}.  Finally,
$\bar\rho_k\ge\bar\rho_0$ gives
\eqref{eq:uniform-x-block-smoothness}.
\end{proof}

The next lemma gives a quasi-descent estimate for the potential function under
the stepsize condition.
\begin{lemma}
\label{lem:descent-skeleton}
\color{blue}
Suppose that Assumptions~\ref{ass:smoothness} and~\ref{ass:h} hold and that
the conditions of Lemma~\ref{lem:x-block-smoothness} are satisfied.  Let
$\theta_x>0$ satisfy $\theta_x\widehat L\le1/2$, and choose
\begin{equation}
    \eta_k=\frac{\theta_x}{\bar\rho_k},
    \qquad k\ge0.
    \label{eq:explicit-primal-stepsize-coupling}
\end{equation}
Then, for every $k\ge0$, it holds that
\begin{equation}\label{eq:variable-penalty-descent-multiplier}
\begin{aligned}
    &\Phi_{s_{k+1}}(\bx^{k+1},\by^{k+1};\bmu^{k+1})
    -\Phi_{s_k}(\bx^k,\by^k;\bmu^k)\\
    &\qquad\le
    -\frac{3\bar\rho_k}{4\theta_x}\|\Delta_{\bx}^{k+1}\|^2
    -\frac{\bar\rho_k}{2}\|\Delta_{\by}^{k+1}\|^2
    +\left(\frac{1}{\tau_k}+\frac{\bar\rho_{k+1}-\bar\rho_k}{2\tau_k^2}\right)
    \|\Delta_{\bmu}^{k+1}\|^2.
\end{aligned}
\end{equation}
\end{lemma}

\begin{proof}
We decompose the one-step change into two parts: the primal descent produced by
the $(\bx,\by)$-updates, and the dual ascent produced by the $\bmu$-update
together with the penalty-parameter drift.

\textbf{For the primal descent.}
\blue{Since $\calS$ is convex and contains $\bx^k$ and $\bx^{k+1}$, the
estimate in Lemma~\ref{lem:x-block-smoothness} applies along the segment
joining these two points.}  Moreover,
\[
    \nabla_{\bx}\Phi_{s_k}(\bx^k,\by^k;\bmu^k)
    =
    \nabla f(\bx^k)+J_{\calA}(\bx^k)^\top\bw^k,
\]
where $\bw^k:=\bmu^k+\bar\rho_k\br^k$.
The $\bx$-update \eqref{eq:x-update-det} gives the identity
\[
    \left\langle \nabla_{\bx}\Phi_{s_k}(\bx^k,\by^k;\bmu^k),\Delta_{\bx}^{k+1}\right\rangle
    =
    -\frac{1}{\eta_k}\|\Delta_{\bx}^{k+1}\|^2,
\]
\blue{so the integral form of the descent lemma gives}
\[
    \Phi_{s_k}(\bx^{k+1},\by^k;\bmu^k)
    -\Phi_{s_k}(\bx^k,\by^k;\bmu^k)
    \le
    -\left(\frac{1}{\eta_k}-\frac{L_k}{2}\right)
    \|\Delta_{\bx}^{k+1}\|^2.
\]
Since $h$ is convex, the exact $\by$-minimization \eqref{eq:y-update} gives
\[
    \Phi_{s_k}(\bx^{k+1},\by^{k+1};\bmu^k)
    -
    \Phi_{s_k}(\bx^{k+1},\by^k;\bmu^k)
    \le
    -\frac{\bar\rho_k}{2}\|\Delta_{\by}^{k+1}\|^2.
\]
Both steps hold $(\bmu^k,\bar\rho_k)$ fixed, so adding them yields the primal
bound
\[
    \Phi_{s_k}(\bx^{k+1},\by^{k+1};\bmu^k)
    -\Phi_{s_k}(\bx^k,\by^k;\bmu^k)
    \le
    -\left(\frac{1}{\eta_k}-\frac{L_k}{2}\right)
    \|\Delta_{\bx}^{k+1}\|^2
    -\frac{\bar\rho_k}{2}\|\Delta_{\by}^{k+1}\|^2.
\]

\textbf{For the dual ascent.}
Holding $(\bx^{k+1},\by^{k+1})$ fixed, the $\bmu$-update and the change of the
penalty coefficient from $\bar\rho_k$ to $\bar\rho_{k+1}$ give
\[
\begin{aligned}
    &\Phi_{s_{k+1}}(\bx^{k+1},\by^{k+1};\bmu^{k+1})
    -\Phi_{s_k}(\bx^{k+1},\by^{k+1};\bmu^k)\\
    &\quad=
    \langle \Delta_{\bmu}^{k+1},\br^{k+1}\rangle
    +\frac{\bar\rho_{k+1}-\bar\rho_k}{2}\|\br^{k+1}\|^2\\
    &\quad=
    \tau_k\|\br^{k+1}\|^2
    +\frac{\bar\rho_{k+1}-\bar\rho_k}{2}\|\br^{k+1}\|^2\\
    &\quad=
    \left(\frac{1}{\tau_k}+\frac{\bar\rho_{k+1}-\bar\rho_k}{2\tau_k^2}\right)
    \|\Delta_{\bmu}^{k+1}\|^2,
\end{aligned}
\]
where the second equality uses \eqref{eq:mu-drift-tau}, i.e.\
$\Delta_{\bmu}^{k+1}=\tau_k\br^{k+1}$, and the last equality substitutes
$\br^{k+1}=\Delta_{\bmu}^{k+1}/\tau_k$.

Combining the two parts yields
\begin{equation}\label{eq:variable-penalty-descent-general}
\begin{aligned}
    &\Phi_{s_{k+1}}(\bx^{k+1},\by^{k+1};\bmu^{k+1})
    -\Phi_{s_k}(\bx^k,\by^k;\bmu^k)\\
    &\qquad\le
    -\left(\frac{1}{\eta_k}-\frac{L_k}{2}\right)
    \|\Delta_{\bx}^{k+1}\|^2
    -\frac{\bar\rho_k}{2}\|\Delta_{\by}^{k+1}\|^2\\
    &\qquad\quad
    +\left(\frac{1}{\tau_k}
    +\frac{\bar\rho_{k+1}-\bar\rho_k}{2\tau_k^2}\right)
    \|\Delta_{\bmu}^{k+1}\|^2.
\end{aligned}
\end{equation}
Moreover,
\begin{equation}
    \frac{1}{\eta_k}-\frac{L_k}{2}
    \ge\frac{3}{4\eta_k}
    =\frac{3\bar\rho_k}{4\theta_x}.
    \label{eq:explicit-primal-descent-coefficient}
\end{equation}
The inequality follows from \eqref{eq:uniform-x-block-smoothness},
\eqref{eq:explicit-primal-stepsize-coupling}, and
$\theta_x\widehat L\le1/2$.  Substituting it into
\eqref{eq:variable-penalty-descent-general} proves
\eqref{eq:variable-penalty-descent-multiplier}.
\end{proof}

\begin{proposition}[Potential upper bound and boundedness consequences]
\label{prop:bounded-iterates}
\color{blue}
Suppose that Assumption~\ref{ass:level-boundedness} and the conditions of
Lemma~\ref{lem:descent-skeleton} hold.  Assume that
\begin{equation}
\begin{aligned}
    C_e
    :=C_r^2\sum_{k=0}^{\infty}
    \left(
        \frac{\tau_k}{\bar\rho_k^2}
        +\frac{\bar\rho_{k+1}-\bar\rho_k}{2\bar\rho_k^2}
    \right)
    <\infty.
\end{aligned}
\label{eq:total-potential-error}
\end{equation}
Define
\begin{equation}
    \overline\Phi
    :=\Phi_{s_0}(\bx^0,\by^0;\bmu^0)+C_e,
    \qquad
    \alpha
    :=\overline\Phi+\frac{C_r^2}{2\bar\rho_0}.
    \label{eq:localized-level-value}
\end{equation}
Then, for every $k\ge0$,
\begin{equation}
    \Phi_{s_k}(\bx^k,\by^k;\bmu^k)\le\overline\Phi,
    \qquad
    F(\bx^k)\le\alpha.
    \label{eq:potential-and-objective-upper-bounds}
\end{equation}
Consequently, define
\begin{equation}
    S:=\calS_\alpha:=\{\bx\in\RR^n:F(\bx)\le\alpha\}.
    \label{eq:derived-x-working-set}
\end{equation}
Then $S$ is compact and
\begin{equation}
    \bx^k\in S,
    \qquad k\ge0,
    \label{eq:derived-x-trajectory-bound}
\end{equation}
Moreover, Lemma~\ref{lem:x-block-smoothness} gives
$\|\by^k\|\le B_y$ for every $k\ge0$.  Thus $\{\bx^k\}$, $\{\by^k\}$,
and $\{\bmu^k\}$ are bounded.  For later use, the constants
\begin{equation}
    G_f:=\sup_{\bx\in S}\|\nabla f(\bx)\|,
    \qquad
    C_A:=\sup_{\bx\in S}\|\calA(\bx)\|
    \label{eq:derived-operator-bounds}
\end{equation}
are finite.
\end{proposition}

\begin{proof}
\color{blue}
Applying \eqref{eq:basic-bounds} to the last term in
\eqref{eq:variable-penalty-descent-multiplier}, dropping the nonpositive
primal terms, and summing from $j=0$ to $k-1$ give
\[
    \Phi_{s_k}(\bx^k,\by^k;\bmu^k)
    \le
    \Phi_{s_0}(\bx^0,\by^0;\bmu^0)+C_e
    =\overline\Phi.
\]
Let $\br^k:=\calA(\bx^k)-\by^k$.  The global Lipschitz continuity of $h$
and the multiplier bound $\|\bmu^k\|\le M_\mu$ yield
\[
\begin{aligned}
    F(\bx^k)
    &\le f(\bx^k)+h(\by^k)+L_h\|\br^k\|\\
    &=\Phi_{s_k}(\bx^k,\by^k;\bmu^k)
      -\langle\bmu^k,\br^k\rangle
      -\frac{\bar\rho_k}{2}\|\br^k\|^2
      +L_h\|\br^k\|\\
    &\le\Phi_{s_k}(\bx^k,\by^k;\bmu^k)
      +C_r\|\br^k\|
      -\frac{\bar\rho_k}{2}\|\br^k\|^2\\
    &\le\overline\Phi+\frac{C_r^2}{2\bar\rho_k}
    \le\alpha.
\end{aligned}
\]
Assumption~\ref{ass:level-boundedness} proves that $\calS_\alpha$ is compact,
and \eqref{eq:potential-and-objective-upper-bounds} gives
$\bx^k\in\calS_\alpha=S$, which proves
\eqref{eq:derived-x-trajectory-bound}.  The bound on $\{\by^k\}$ follows from
Lemma~\ref{lem:x-block-smoothness}, and boundedness of $\{\bmu^k\}$ follows
from Lemma~\ref{lem:multiplier-boundedness}.  Finally,
continuity of $\nabla f$ and $\calA$ on the compact set $S$ proves
\eqref{eq:derived-operator-bounds}.
\end{proof}

\begin{proposition}[Weighted finite length]
\label{lem:weighted-finite-length}
\blue{Suppose that Assumption~\ref{ass:level-boundedness} and the conditions of
Lemma~\ref{lem:descent-skeleton} hold.}
Assume in addition that
\[
    \bar\rho_k=\rho+c_s(k+1)^a,
    \qquad
    c_s>0,\quad a\in(0,1).
\]
Assume also that the dual stepsizes satisfy
\begin{equation}
    \sum_{k=0}^{\infty}
    \frac{\tau_k}{\bar\rho_k^2}
    <\infty.
    \label{eq:dual-ascent-summability}
\end{equation}
Then there exists a constant $C_{\mathrm{desc}}>0$, independent of $K$,
such that, for every $K\ge1$,
\begin{equation}\label{eq:weighted-finite-length}
    \sum_{k=0}^{K-1}
    \bar\rho_k\left(\|\Delta_{\bx}^{k+1}\|^2+\|\Delta_{\by}^{k+1}\|^2\right)
    \le
    C_{\mathrm{desc}} .
\end{equation}
\end{proposition}

\begin{proof}
Throughout the proof, $C>0$ denotes a generic constant independent of $k$
and $K$.  The value of $C$ may change from line to line and depends only on
the problem data and fixed algorithmic parameters.

\textbf{Step 1: Summability of the controlled ascent error.}
Define
\[
    e_k
    :=
    C_r^2\left(
    \frac{\tau_k}{\bar\rho_k^2}
    +
    \frac{\bar\rho_{k+1}-\bar\rho_k}{2\bar\rho_k^2}
    \right).
\]
The first series in $e_k$ is finite by
\eqref{eq:dual-ascent-summability}.  For the penalty-drift series, the
concavity of $t\mapsto t^a$ and
$\bar\rho_k=\rho+c_s(k+1)^a$ give
\[
    \bar\rho_{k+1}-\bar\rho_k
    =
    c_s\bigl((k+2)^a-(k+1)^a\bigr)
    \le
    c_sa(k+1)^{a-1}.
\]
Moreover, $\bar\rho_k\ge c_s(k+1)^a$, so
\[
    \frac{\bar\rho_{k+1}-\bar\rho_k}{\bar\rho_k^2}
    \le
    \frac{c_sa(k+1)^{a-1}}{c_s^2(k+1)^{2a}}
    =
    O(k^{-(1+a)}).
\]
Consequently, $\sum_{k=0}^{\infty}e_k<\infty$.
\blue{Proposition~\ref{prop:bounded-iterates} now applies.  Hence the compact
set $S=\calS_\alpha$ defined in
\eqref{eq:derived-x-working-set} contains every $\bx^k$.}

\textbf{Step 2: Telescoping descent.}
In \eqref{eq:variable-penalty-descent-multiplier}, the ascent term
$\bigl(\frac{1}{\tau_k}+\frac{\bar\rho_{k+1}-\bar\rho_k}{2\tau_k^2}\bigr)
\|\Delta_{\bmu}^{k+1}\|^2$ rewrites as
$\tau_k\|\br^{k+1}\|^2+\frac{\bar\rho_{k+1}-\bar\rho_k}{2}\|\br^{k+1}\|^2$
via $\Delta_{\bmu}^{k+1}=\tau_k\br^{k+1}$, and is bounded by $e_k$ using
\eqref{eq:basic-bounds}.
By \eqref{eq:variable-penalty-descent-multiplier}, \blue{set
$c_{\mathrm{desc}}:=\min\{3/(4\theta_x),1/2\}>0$.  Then}
\begin{equation}\label{eq:one-step-descent-clean}
    \Phi_{s_{k+1}}(\bx^{k+1},\by^{k+1};\bmu^{k+1})
    -
    \Phi_{s_k}(\bx^k,\by^k;\bmu^k)
    \le
    -c_{\mathrm{desc}}\bar\rho_k
    \left(\|\Delta_{\bx}^{k+1}\|^2+\|\Delta_{\by}^{k+1}\|^2\right)
    +e_k,
\end{equation}
Summing \eqref{eq:one-step-descent-clean} from $k=0$ to $K-1$ yields
\[
    c_{\mathrm{desc}}
    \sum_{k=0}^{K-1}
    \bar\rho_k\left(\|\Delta_{\bx}^{k+1}\|^2+\|\Delta_{\by}^{k+1}\|^2\right)
    \le
    \Phi_{s_0}(\bx^0,\by^0;\bmu^0)
    -
    \Phi_{s_K}(\bx^K,\by^K;\bmu^K)
    +
    \sum_{k=0}^{K-1}e_k.
\]

\textbf{Step 3: Lower bound for the potential.}
{\color{blue}
Since $h$ is globally $L_h$-Lipschitz and
$\br^K=\calA(\bx^K)-\by^K$,
\[
    h(\by^K)
    \ge h\bigl(\calA(\bx^K)\bigr)-L_h\|\br^K\|.
\]
Using this inequality, $\|\bmu^K\|\le M_\mu$, and completing the square give
\[
\begin{aligned}
    \Phi_{s_K}(\bx^K,\by^K;\bmu^K)
    &=
    f(\bx^K)+h(\by^K)
    +\langle\bmu^K,\br^K\rangle
    +\frac{\bar\rho_K}{2}\|\br^K\|^2\\
    &\ge
    F(\bx^K)-C_r\|\br^K\|
    +\frac{\bar\rho_K}{2}\|\br^K\|^2\\
    &\ge
    \inf_{\bx\in S}F(\bx)
    -\frac{C_r^2}{2\bar\rho_K}\\
    &\ge
    \inf_{\bx\in S}F(\bx)
    -\frac{C_r^2}{2\bar\rho_0}.
\end{aligned}
\]
The infimum is finite because $F$ is continuous on the compact set $S$.
Therefore,
\[
    \sum_{k=0}^{K-1}
    \bar\rho_k\left(\|\Delta_{\bx}^{k+1}\|^2+\|\Delta_{\by}^{k+1}\|^2\right)
    \le
    \frac{1}{c_{\mathrm{desc}}}
    \left[
    \Phi_{s_0}(\bx^0,\by^0;\bmu^0)
    -
    \inf_{\bx\in S}F(\bx)
    +
    \frac{C_r^2}{2\bar\rho_0}
    +
    \sum_{k=0}^{\infty}e_k
    \right]
    =:
    C_{\mathrm{desc}},
\]
where $C_{\mathrm{desc}}$ is finite and independent of $K$.
}
\end{proof}

\begin{remark}[Dual stepsize design]
\label{rmk:dual-stepsize-choice}
Lemma~\ref{lem:multiplier-boundedness} shows that multiplier boundedness only
requires $0<\tau_k\le s_k$.  \blue{The additional condition needed for the
descent argument is}
\[
    \sum_{k=0}^{\infty}\frac{\tau_k}{\bar\rho_k^2}
    <\infty.
\]
For $s_k=c_s(k+1)^a$, a family satisfying
$\tau_k=O\bigl(s_k/(k+1)^q\bigr)$ meets this requirement whenever
$q>1-a$.  If a downstream result additionally requires finite total variation
of the anchor multiplier, Lemma~\ref{lem:multiplier-boundedness} shows that the
stronger condition $\sum_k\tau_k/\bar\rho_k<\infty$ is sufficient.  The global
complexity result below uses $q=1$, namely
$\tau_k=\theta_\mu s_k/(k+2)$.  This choice makes the potential-ascent error
summable and gives
$\operatorname{dist}(\bw^k,\partial h(\by^k))=O(k^{-1})$; finite total
variation of $\{\bmu^k\}$ is not needed for that result.
\end{remark}
\begin{theorem}[Global KKT residual complexity]
\label{theorem:target-complexity}
\blue{Suppose that Assumptions~\ref{ass:level-boundedness},
\ref{ass:smoothness}, and~\ref{ass:h} hold.}  Let $c_s>0$, $a\in(0,1)$,
$\theta_x>0$ with $\theta_x\widehat L\le1/2$, and
$\theta_\mu\in(0,1]$, and choose
\begin{equation}
\label{eq:complexity-parameter-schedule}
    s_k=c_s(k+1)^a,
    \qquad
    \bar\rho_k=\rho+s_k,
    \qquad
    \eta_k=\frac{\theta_x}{\bar\rho_k},
    \qquad
    \tau_k=\frac{\theta_\mu s_k}{k+2},
\end{equation}
where $\widehat L$ is the constant in
\eqref{eq:uniform-x-block-smoothness}.

For any $K\ge2$, choose
\[
    \kappa
    \in
    \argmin_{k\in\{\lceil K/2\rceil,\ldots,K-1\}}
    \bar\rho_k\left(\|\Delta_{\bx}^{k+1}\|^2+\|\Delta_{\by}^{k+1}\|^2\right).
\]
Define the effective KKT multiplier
\[
    \bz^\kappa
    :=
    \bmu^{\kappa-1}+\bar\rho_{\kappa-1}\br^\kappa,
\]
where $\br^\kappa:=\calA(\bx^\kappa)-\by^\kappa$.
Then
\begin{equation}\label{eq:kkt-residual-general-a}
\begin{aligned}
    \left\|
        \nabla f(\bx^\kappa)
        +J_{\calA}(\bx^\kappa)^\top \bz^\kappa
    \right\|
    &=O\!\left(K^{-(1-a)/2}\right),\\
    \operatorname{dist}\bigl(\bz^\kappa,\partial h(\by^\kappa)\bigr)
    &=0,\\
    \|\br^\kappa\|
    &=O(K^{-a}).
\end{aligned}
\end{equation}
Moreover, $\|\bz^\kappa\|\le L_h$.  In particular, for the balanced
choice $a=1/3$,
\[
    \max\left\{
        \left\|
            \nabla f(\bx^\kappa)
            +J_{\calA}(\bx^\kappa)^\top \bz^\kappa
        \right\|,
        \operatorname{dist}\bigl(
            \bz^\kappa,\partial h(\by^\kappa)
        \bigr),
        \|\br^\kappa\|
    \right\}
    =O(K^{-1/3}).
\]
Consequently, to obtain an $\epsilon$-approximate KKT point in the sense that
\[
    \left\|
    \nabla f(\bx^\kappa)
    +
    J_{\calA}(\bx^\kappa)^\top \bz^\kappa
    \right\|
    \le
    \epsilon,
    \qquad
    \operatorname{dist}\bigl(
        \bz^\kappa,\partial h(\by^\kappa)
    \bigr)
    \le
    \epsilon,
    \qquad
    \|\br^\kappa\|
    \le
    \epsilon,
\]
under the assumptions above, it suffices to take
\[
    K
    =
    O(\epsilon^{-3}).
\]
Moreover, the virtual multiplier
$\blmbd^\kappa=\bmu^{\kappa-1}+s_{\kappa-1}\br^\kappa$
and the KKT-aligned multiplier
$\bz^\kappa=\bmu^{\kappa-1}+\bar\rho_{\kappa-1}\br^\kappa$
satisfy
\[
    \bz^\kappa-\blmbd^\kappa
    =
    \rho\br^\kappa,
\]
and hence
\[
    \operatorname{dist}\bigl(
        \blmbd^\kappa,\partial h(\by^\kappa)
    \bigr)
    \le \rho\|\br^\kappa\|
    =O(K^{-a}).
\]
The multiplier driving the $\bx$-step also satisfies
\[
    \operatorname{dist}\bigl(
        \bw^\kappa,\partial h(\by^\kappa)
    \bigr)
    =O(K^{-1}).
\]
The KKT residual is stated with $\bz^\kappa$ because
$\bz^\kappa\in\partial h(\by^\kappa)$ by the preceding $\by$-step
optimality condition.
\end{theorem}

\begin{proof}
The primal stepsize in \eqref{eq:complexity-parameter-schedule} satisfies
\eqref{eq:explicit-primal-stepsize-coupling} and the conditions of
Lemma~\ref{lem:descent-skeleton}.
Moreover, $0<\tau_k\le s_k$ and
\[
\begin{aligned}
    \sum_{k=0}^{\infty}\frac{\tau_k}{\bar\rho_k^2}
    &\le
    \theta_\mu
    \sum_{k=0}^{\infty}
    \frac{1}{(k+2)\bar\rho_k}\\
    &\le
    \frac{\theta_\mu}{c_s}
    \sum_{k=0}^{\infty}
    \frac{1}{(k+2)(k+1)^a}
    <\infty.
\end{aligned}
\]
Thus all conditions of Proposition~\ref{lem:weighted-finite-length} hold, and
that proposition gives
\[
    \sum_{k=0}^{K-1}
    \bar\rho_k\left(\|\Delta_{\bx}^{k+1}\|^2+\|\Delta_{\by}^{k+1}\|^2\right)
    \le
    C_{\mathrm{desc}},
    \qquad
    \forall K\ge1.
\]
The index set $\{\lceil K/2\rceil,\ldots,K-1\}$ has
$\lfloor K/2\rfloor$ elements, and $\lfloor K/2\rfloor\ge K/3$ for every
$K\ge2$.  Hence, by the definition of $\kappa$,
\[
    \bar\rho_\kappa\left(\|\Delta_{\bx}^{\kappa+1}\|^2+\|\Delta_{\by}^{\kappa+1}\|^2\right)
    \le
    \frac{3C_{\mathrm{desc}}}{K}.
\]
Since $\kappa\in\{\lceil K/2\rceil,\ldots,K-1\}$ and
$\bar\rho_k=\rho+c_s(k+1)^a$, we have $\bar\rho_\kappa\asymp K^a$.  Thus
\[
    \|\Delta_{\bx}^{\kappa+1}\|+\|\Delta_{\by}^{\kappa+1}\|
    =
    O\!\left(K^{-(1+a)/2}\right).
\]

\textbf{Step 1: Feasibility and multiplier residuals.}
By Lemma~\ref{lem:multiplier-boundedness},
\[
    \|\br^\kappa\|
    \le
    \frac{C_r}{\bar\rho_{\kappa-1}}
    =
    O(K^{-a}).
\]
The exact $\by$-update optimality condition and
\eqref{eq:bounded-h-subgradients} give
\[
    \bz^\kappa\in\partial h(\by^\kappa),
    \qquad
    \|\bz^\kappa\|\le L_h,
\]
which proves the second residual in
\eqref{eq:kkt-residual-general-a}.  Moreover,
$\bz^\kappa-\blmbd^\kappa=\rho\br^\kappa$ gives the stated virtual-multiplier
residual.

For every $k\ge1$, the definitions of $\bw^k$ and $\bz^k$, together with
\eqref{eq:mu-drift-tau}, give
\[
\begin{aligned}
    \bw^k-\bz^k
    &=
    \bmu^k-\bmu^{k-1}
    +\bigl(\bar\rho_k-\bar\rho_{k-1}\bigr)\br^k\\
    &=
    \bigl(
        \tau_{k-1}+\bar\rho_k-\bar\rho_{k-1}
    \bigr)\br^k.
\end{aligned}
\]
Using \eqref{eq:basic-bounds} with index $k-1$ and the schedule
\eqref{eq:complexity-parameter-schedule}, we obtain
\[
\begin{aligned}
    \|\bw^k-\bz^k\|
    &\le
    C_r\left(
        \frac{\tau_{k-1}}{\bar\rho_{k-1}}
        +
        \frac{\bar\rho_k-\bar\rho_{k-1}}
             {\bar\rho_{k-1}}
    \right)\\
    &\le
    C_r\left(
        \frac{\theta_\mu}{k+1}
        +\frac{a}{k}
    \right)
    =O(k^{-1}),
\end{aligned}
\]
where the second inequality uses
$(k+1)^a-k^a\le ak^{a-1}$.  Since the exact $\by$-update optimality
condition gives $\bz^k\in\partial h(\by^k)$,
\[
    \operatorname{dist}\bigl(\bw^k,\partial h(\by^k)\bigr)
    \le \|\bw^k-\bz^k\|
    =O(k^{-1}).
\]
In particular, $\kappa\ge K/2$ gives
$\|\bw^\kappa-\bz^\kappa\|=O(K^{-1})$.

\textbf{Step 2: Stationarity residual.}
The $\bx$-update optimality condition gives
\[
    \nabla f(\bx^\kappa)
    +
    J_{\calA}(\bx^\kappa)^\top \bw^\kappa
    =
    -\frac{1}{\eta_\kappa}\Delta_{\bx}^{\kappa+1},
\]
where $\bw^\kappa:=\bmu^\kappa+\bar\rho_\kappa\br^\kappa$.
Consequently,
\[
\begin{aligned}
    \left\|
        \nabla f(\bx^\kappa)
        +J_{\calA}(\bx^\kappa)^\top \bz^\kappa
    \right\|
    &\le
    \frac{1}{\eta_\kappa}\|\Delta_{\bx}^{\kappa+1}\|
    +G_A\|\bw^\kappa-\bz^\kappa\|\\
    &=O\!\left(K^{-(1-a)/2}\right)+O(K^{-1})\\
    &=O\!\left(K^{-(1-a)/2}\right).
\end{aligned}
\]

The two residual rates are balanced by solving
$a=(1-a)/2$, which gives $a=1/3$. For this choice, both stationarity and
splitting residuals are $O(K^{-1/3})$, and the iteration complexity claim follows by
setting $K=O(\epsilon^{-3})$.
\end{proof}

\begin{remark}
\label{rmk:first-order-eps-minus-three}
The schedule \eqref{eq:complexity-parameter-schedule} couples the dual
regularization, primal step, and dual step as
\[
    \frac{1}{s_k}=\Theta(k^{-a}),
    \qquad
    \eta_k=\Theta(k^{-a}),
    \qquad
    \tau_k=\Theta(k^{a-1}).
\]
The dual-ascent and penalty-drift errors then obey
\[
    \frac{\tau_k}{\bar\rho_k^2}
    =O(k^{-(1+a)}),
    \qquad
    \frac{\bar\rho_{k+1}-\bar\rho_k}{\bar\rho_k^2}
    =O(k^{-(1+a)}),
\]
so both error sequences are summable.  In addition,
\[
    \|\Delta_{\bmu}^{k+1}\|=O(k^{-1}),
    \qquad
    \operatorname{dist}\bigl(\bw^k,\partial h(\by^k)\bigr)=O(k^{-1}).
\]
Thus the splitting residual is retained as the feasibility component of the
KKT residual, while its contribution to the potential ascent is converted to
the summable multiplier-increment term
\[
    \tau_k\|\br^{k+1}\|^2
    =
    \frac{1}{\tau_k}\|\Delta_{\bmu}^{k+1}\|^2.
\]
Finally, weighted finite length gives
\[
    r_{\rm stat}=O(K^{-(1-a)/2}),
    \qquad
    r_{\rm feas}=O(K^{-a}).
\]
Balancing these rates yields $a=1/3$.  Hence
\[
    \frac{1}{s_k}=\Theta(k^{-1/3}),
    \qquad
    \eta_k=\Theta(k^{-1/3}),
    \qquad
    \tau_k=\Theta(k^{-2/3}),
\]
and all three KKT residuals are $O(K^{-1/3})$.
\end{remark}
\subsection{A conditional local fast rate}
\label{subsec:local-fast-rate}

The preceding theorem is deliberately stated without local or global error
bound assumptions.  If one is willing to impose a stronger local condition
that rules out nearby infeasible stationary points of the reduced augmented
Lagrangian, the same algorithmic template admits a faster local rate with a
constant effective penalty.  This statement is conditional and should be read
as a local refinement rather than as a replacement for
Theorem~\ref{theorem:target-complexity}.

\begin{assumption}[Local reduced error bound]
\label{ass:local-reduced-eb}
There exist an index $k_0\ge0$, a constant $\kappa_{\rm eb}>0$, and a
neighborhood $\mathcal U$ of the KKT solution set such that
$(\bx^k,\bmu^k)\in\mathcal U$ for all $k\ge k_0$, and
\begin{equation}\label{eq:local-reduced-eb}
    \|\calA(\bx^k)\|
    \le
    \kappa_{\rm eb}
    \left\|
    \nabla f(\bx^k)
    +
    J_{\calA}(\bx^k)^\top
    \bigl(\bmu^k+\bar\rho\,\calA(\bx^k)\bigr)
    \right\|,
    \qquad
    \forall k\ge k_0 .
\end{equation}
\end{assumption}

\begin{remark}
Assumption~\ref{ass:local-reduced-eb} is stronger than uniform LICQ.
Uniform LICQ controls multiplier size from a stationarity residual, whereas
\eqref{eq:local-reduced-eb} also requires the feasibility violation to be
controlled by the stationarity residual of the fixed-penalty reduced model.
Equivalently, it excludes local infeasible stationary points of
$\Phi(\cdot;\bmu^k)$ along the generated anchor sequence.
\end{remark}

\begin{corollary}[Local fast rate under a reduced error bound]
\label{cor:local-fast-rate}
\blue{Suppose that the conditions of Proposition~\ref{prop:bounded-iterates},
Assumption~\ref{ass:strong-licq}, and
Assumption~\ref{ass:local-reduced-eb} hold.}  Let the effective penalty be
constant, namely $\bar\rho_k\equiv\bar\rho>0$, and choose a constant primal
stepsize $\eta>0$ such that
\[
    \frac{1}{2\eta}
    -
    \frac{1}{2}
    \bigl(
    L_f+L_AM_\mu+\bar\rho(L_AC_A+G_A^2)
    \bigr)
    \ge
    \alpha
\]
for some $\alpha>0$.
Assume also that
\begin{equation}
    \sum_{k=0}^{\infty}\frac{\tau_k}{\bar\rho}<\infty.
    \label{eq:anchor-total-variation}
\end{equation}
For any $K>k_0$, choose
\[
    \kappa
    \in
    \argmin_{k\in\{k_0,\ldots,K-1\}}
    \|\Delta_{\bx}^{k+1}\|^2,
\]
and define
\[
    \by^\kappa
    :=
    \bmu^\kappa+\bar\rho\,\calA(\bx^\kappa).
\]
Then there exists a constant $C_{\rm fast}>0$, independent of $K$, such that
\[
    \left\|
    \nabla f(\bx^\kappa)
    +
    J_{\calA}(\bx^\kappa)^\top \by^\kappa
    \right\|^2
    +
    \|\calA(\bx^\kappa)\|^2
    \le
    \frac{C_{\rm fast}}{K-k_0}.
\]
Consequently, a point satisfying both
\[
    \left\|
    \nabla f(\bx^\kappa)
    +
    J_{\calA}(\bx^\kappa)^\top \by^\kappa
    \right\|
    \le
    \epsilon,
    \qquad
    \|\calA(\bx^\kappa)\|\le\epsilon
\]
is obtained after $K-k_0=O(\epsilon^{-2})$ iterations.
\end{corollary}

\begin{proof}
With $\bar\rho_k\equiv\bar\rho$, the penalty-drift term in
\eqref{eq:variable-penalty-descent-general} vanishes.  The stated
stepsize condition gives
\[
    \Phi(\bx^{k+1};\bmu^{k+1})-\Phi(\bx^k;\bmu^k)
    \le
    -\alpha\|\Delta_{\bx}^{k+1}\|^2
    +
    C_A\|\Delta_{\bmu}^{k+1}\|,
    \qquad
    \forall k\ge0,
\]
where
\[
    \Phi(\bx;\bmu)
    :=
    f(\bx)+\langle\bmu,\calA(\bx)\rangle
    +\frac{\bar\rho}{2}\|\calA(\bx)\|^2 .
\]
Summing from $k=k_0$ to $K-1$, using the lower boundedness of $\Phi$ on the
compact working set and the summability of
$\{\|\Delta_{\bmu}^{k+1}\|\}$ implied by
\eqref{eq:basic-bounds} and
\eqref{eq:anchor-total-variation},
yields
\[
    \sum_{k=k_0}^{K-1}\|\Delta_{\bx}^{k+1}\|^2
    \le
    C
\]
for a constant $C$ independent of $K$.  Hence the definition of $\kappa$
implies
\[
    \|\Delta_{\bx}^{\kappa+1}\|^2
    \le
    \frac{C}{K-k_0}.
\]
The explicit primal step gives
\[
    \nabla f(\bx^\kappa)
    +
    J_{\calA}(\bx^\kappa)^\top \by^\kappa
    =
    \frac{1}{\eta}(\bx^\kappa-\bx^{\kappa+1}),
\]
so the squared stationarity residual is bounded by
$C/(\eta^2(K-k_0))$.  Assumption~\ref{ass:local-reduced-eb} then gives
\[
    \|\calA(\bx^\kappa)\|^2
    \le
    \kappa_{\rm eb}^2
    \left\|
    \nabla f(\bx^\kappa)
    +
    J_{\calA}(\bx^\kappa)^\top \by^\kappa
    \right\|^2
    \le
    \frac{C}{K-k_0}.
\]
Combining the two estimates proves the claim.
\end{proof}
